% ============================================================
% 建模与求解 — 内部报告子文件模板
% 包含：模型建立（三段式）+ 模型求解
% 遵循 PR-015：模型建立三段式与编号步骤法
% 依赖包：float, longtable, array
% ============================================================

\subsubsection{模型建立}

% ==== 第1段：开篇引文 ====
针对问题一，本节将从...角度建立...模型。具体过程如下。

% ==== 第2段：编号步骤法（每步≥6公式+文字串联+变量说明）====
\textbf{步骤1：...}

首先，...。设...为...，则有：
\begin{equation}
    ... = ...
\end{equation}
其中，$...$表示...。

将...代入...，可得：
\begin{equation}
    ... = ...
\end{equation}
进一步展开得：
\begin{equation}
    ... = ...
\end{equation}

\textbf{步骤2：...}

根据...的性质，可以建立以下关系：
\begin{equation}
    ... = ...
\end{equation}
其中，$...$为...，$...$为...。

由上述公式联立可得：
\begin{equation}
    ... = ...
\end{equation}

% ...（按需增加步骤，每步≥6公式）...

\textbf{步骤N：求解算法}

% 算法参数表（longtable 格式，PR-013 列宽公式：4列合计 0.88\textwidth）
\begin{longtable}{>{\centering\arraybackslash}p{0.12\textwidth}>{\centering\arraybackslash}p{0.30\textwidth}>{\centering\arraybackslash}p{0.20\textwidth}>{\centering\arraybackslash}p{0.26\textwidth}}
  \caption{算法参数设置}
  \label{tab:参数1} \\
  \toprule
  参数名 & 参数含义 & 取值 & 选取依据 \\
  \midrule
  \endfirsthead
  \caption{算法参数设置（续）} \\
  \toprule
  参数名 & 参数含义 & 取值 & 选取依据 \\
  \midrule
  \endhead
  \bottomrule
  \endfoot
  ... & ... & ... & ... \\
\end{longtable}

% ==== 第3段：总结桥接 ====
综上所述，本节完成了...模型的建立，得到了...（目标函数/评价体系/预测模型），接下来将基于此进行求解。

% ============================================================
\subsubsection{模型求解}

% 引文
基于上述模型，本节利用...数据进行求解。核心参数如表\ref{tab:参数1}所示。

% 结果表格（longtable，PR-013 列宽公式：5列合计 0.84\textwidth）
求解得到以下结果：

\begin{longtable}{>{\centering\arraybackslash}p{0.12\textwidth}>{\centering\arraybackslash}p{0.18\textwidth}>{\centering\arraybackslash}p{0.18\textwidth}>{\centering\arraybackslash}p{0.18\textwidth}>{\centering\arraybackslash}p{0.18\textwidth}}
  \caption{...结果}
  \label{tab:结果1} \\
  \toprule
  指标1 & 指标2 & 指标3 & 指标4 & 指标5 \\
  \midrule
  \endfirsthead
  \caption{...结果（续）} \\
  \toprule
  指标1 & 指标2 & 指标3 & 指标4 & 指标5 \\
  \midrule
  \endhead
  \bottomrule
  \endfoot
  ... & ... & ... & ... & ... \\
\end{longtable}

% 结果图
如图\ref{fig:xxx}所示，...（一句话点出图的结论）。

% 图引用示例（在求解代码中生成图后，此处引用）：
% \begin{figure}[H]
%   \centering
%   \includegraphics[width=0.8\textwidth]{../求解/问题一/figures/xxx.png}
%   \caption{...}
%   \label{fig:xxx}
% \end{figure}

% 结果分析
由表\ref{tab:结果1}可知，...。从图\ref{fig:xxx}可以看出，...。
综合以上结果，...。
